\chapter{镁\texorpdfstring{\quad}{ }铝}\label{chp:Mg_Al}
在已发现的一百多种元素里，大约有五分之四是金属元素。
金属在自然界里分布很广，无论在矿物、动植物或水中，都可以发现或多或少的金属元素存在。少数化学性质不活泼的金属，在自然界里以单质形态存在，化学性质活泼的金属，总是以化合态存在的。
金属有不同的分类方法 在冶金工业上，人们常把金属分为黑色金属（包括铁、铬、锰）和有色金属（铁、 铬、锰以外的金属）两大类。
根据金属的密度不同，还可以把金属分为重金属和轻金属，密度大于 \qty{4.5}{g/cm^3} 的叫做重金属（如铜、镍、铅等）； 密度小于 \qty{4.5}{g/cm^3} 的叫做轻金属（如钾、钠、钙、镁、铝等）。
除此以外，还可把金属分为常见金属（如铁、铝等）和稀有金属（如锆、铪、铌、钼等）。
这一章主要学习两种重要的金属：镁和铝。
为了更好地了解金属的一些通性，让我们先学习关于金属键的知识。

\section{金属键}
金属有许多共同的性质，象有金属光泽、不透明、容易导电、导热、有延展性，等等。怎样解释金属的这些共同的性质呢？

我们学习离子晶体、分子晶体和原子晶体的时候，已经知道这些晶体都有不同的特性，而这些特性一般是由它们的晶体结构所决定的。
同样，我们可以想象，金属的这些共同的性质也是由金属的结构所决定的。

金属（除汞外）在常温下，一般都是晶体。
如\cref{fig:6-1} 就是铝晶体结构的示意图，用 X 射线进行研究发现，铝原子好象很多硬球一层一层很紧密地堆积，形成了晶体。
而且，每一个铝原子的周围都有较多的铝原子围绕着。
其他金属的结构跟铝相似，它们的结构也都是由金属原子紧密堆积而形成的晶体，每一个金属原子也有许多相同的原子围绕着。
那么，金属原子又是怎样结合在一起而构成金属晶体的呢?

\par\medskip\noindent
\begin{minipage}{0.65\linewidth}\parindent2em
我们知道，金属原子的价电子比较少，电离能又比较低，所以金属原子容易失去电子。
因此，金属的结构实际上是金属原子释出电子后所形成的金属离子按一定规律堆积，释出的价电子在整个晶体里自由地运动着。
这些电子叫做自由电子。
金属离子跟自由电子之间存在着较强的作用，因而使许多金属离子相互结合在一起。
象这种金属离子跟自由电子之间存在的较强的作用，叫做\Concept{金属键}。
我们可以把金属键想象为: “释去了价电子的金属离子沉浸在自由电子的海洋里。”
通过金属键形成的单质晶体，叫做\Concept{金属晶体}。
\end{minipage}\hfill
\begin{minipage}{0.33\linewidth}
\begin{figurehere}
  \includegraphics{6-1.pdf}
  \caption{铝晶体结构示意图}\label{fig:6-1}
\end{figurehere}
\end{minipage}\par\medskip

金属键跟共价键不同，共价键具有方向性，金属键没有方向性。
金属的价电子的数目比较少，而围绕在每一个金属离子周围的金属离子的数目又很多，所以，金属键不象共价键那样能形成共用电子对。
因为一个金属原子的少数价电子不能同时跟周围许多金属原子的价电子形成共用电子对。
所以，金属晶体里的自由电子不专属于某几个特定的金属离子，而为许多金属离子所共有，它们几乎均匀地分布在整个晶体里。

现在让我们用金属键的知识来简单地说明金属的一些共同的性质。

在通常情况下，金属晶体里自由电子的运动是没有一定方向的。
但在外加电场的条件下，自由电子在金属晶体里就会发生定向运动，因而形成电流。
这就是金属容易导电的原因。

金属的导热性也跟金属晶体里自由电子的运动有关。自由电子在运动时经常跟金属离子相碰撞，从而引起两者能量的交换。当金属某一部分受热时，在那个区域里的自由电子的能量增加，运动速度也随之加快，于是通过碰撞，自由电子就能把能量传给其他金属离子。金属就是借着自由电子的运动把能量从温度高的部分传到温度低的部分，从而使整块金属达到同样的温度。

金属的延展性也可以从金属晶体的结构特点加以说明。当金属受外力作用时，各层之间就发生了相对的滑动（\cref{fig:6-2}），但是由于金属键没有方向性，滑动以后，各层之间仍保持金属键的作用，因此，在外力作用下，金属虽然发生了形变，但金属不致断裂，所以金属一般地说都有不同程度的延展性。
\begin{figure}
  \includegraphics{6-2.pdf}
  \caption{金属延展性示意图}\label{fig:6-2}
\end{figure}

上面，我们解释了金属的一些共同的性质，但不同金属之间，在某些性质（如密度、硬度、熔点等）方面，又表现出很大的差别。这些性质与金属原子本身的性质，金属原子的堆积方式等因素有关。

\begin{Practice}[习题]
\begin{question}
  \item 金属键跟共价键有什么区别？
  \item 为什么金属容易导电、导热和有延展性，而在通常情况下，离子化合物却没有上述性质？
  \item 今有五种处于固态的物质: 钠、硅、氛、氯化钠、冰。 问下列五项性质的叙述各适用于哪种物质？
  \begin{tasks}[label=(\arabic*)]
    \task 由分子间力结合而成，熔点很低。
    \task 电的良导体,熔点在 \qty{100}{\celsius} 左右。
    \task 由共价键组成的网状晶体，熔点很高。
    \task 非导体，但熔融后可以导电。
    \task 在晶体结构里有氢键。
  \end{tasks}
  \item 金属导电和电解质溶液导电的方式有什么不同？
  \item 试举例说明离子晶体、分子晶体、原子晶体和金属晶体在结构上、性质上有什么不同。
\end{question}
\end{Practice}

\section{镁和铝的性质}
镁和铝都是第三周期的轻金属元素，镁原子核外电子排布为 $1s^22s^22p^63s^2$，铝原子核外电子排布为 $1s^22s^22p^63s^23p^1$。
这两种原子的价电子数都比较少，因而它们的性质有类似的地方。
但是由于结构不同，所以它们的性质又有许多差异。
为了便于比较，我们把镁和铝的性质并列介绍。

\subsection{物理性质}
\cref{tab:6-1} 是镁和铝的原子和单质的一些性质。
\begin{table}
  \caption{镁和铝的原子和单质的一些性质}\label{tab:6-1}
  \begin{tblr}{
    colspec={c*2{X[2,c]}cccc*3{X[c]}},hline{3}=0.8pt,
  }
    \SetCell[r=2]{m,c}{元素\\符号} & \SetCell[r=2]{m,c} {原子半径 \\ (\qty{e-10}{m})} & \SetCell[r=2]{m,c} {离子半径\\ (\qty{e-10}{m})} & \SetCell[r=2]{m,c} {第一电离能 \\ (\unit{eV}) } & \SetCell[r=2]{m,c} {化\\合\\价}&\SetCell[c=5]{m,c} 单质&&&& \\
    &&&&&{颜\\色}&{密度\\ (\unit{g/cm^3})}&{熔点\\ (\unit{\celsius})}&{沸点\\ (\unit{\celsius})}&{硬度}\\
    \ce{Mg}& 1.60 & 0.66 & 7.644 & $+2$ & {银\\白} & 1.74 & 648.8 & 1090 & 很软\\
    \ce{Al}& 1.43 & 0.51 & 5.984 & $+3$ & {银\\白} & 2.70 & 660.4 & 2467 & 较软\\
  \end{tblr}\par\smallskip
  \begin{minipage}{\linewidth}
  \footnotesize 表中的离子半径是指离子晶体里的 \ce{Al^{3+}} 和 \ce{Mg^{2+}} 的半径。
  \end{minipage}
\end{table}

从\cref{tab:6-1} 可以知道，镁和铝都是密度较小、熔点也比较低、硬度较小的银白色金属。
但镁和铝比较，铝的硬度比镁的稍大，熔点和沸点都比镁的高。
这主要是因为镁和铝的金属键的强弱不同。
我们知道，铝原子的价电子数比镁原子的多，它的核电荷比镁的大，而它的原子半径却比镁的小。
所以在铝的晶体里，自由电子跟铝离子之间的作用比较强。
因此，铝的硬度比镁的大，熔点和沸点都比镁的高。

纯铝的导电性很好，在电力工业上它可以代替部分铜作导线和电缆。
铝有很大的延展性，能够抽成细丝，也能压成薄片成为铝箔，铝箔可以用来包装胶卷、糖果等。
铝粉跟某些油料混和，可以制银白色防锈油漆。
铝的更重要的用途是可以跟许多元素形成合金。
所谓\Concept{合金}就是\emph{两种或两种以上的金属（或金属跟非金属）熔合而成的具有金属特性的物质}。
合金比它的成分金属具有许多良好的物理的、化学的或机械的性能。
一般地说，多数合金的熔点比它的各成分金属的熔点都低。
例如，铝硅合金（含 \ce{Si} 13.5\%）的熔点为 \qty{564}{\celsius}，比纯铝或硅的熔点低，而且它在凝固时收缩率又很小。
因而，这种合金适合铸造。
又如，硬铝（含 \ce{Cu} 4\%、\ce{Mg} 0.5\%、\ce{Mn} 0.5\%、\ce{Si} 0.7\%）的强度和硬度都比纯铝大，几乎相当于钢材，而密度又小。
铝合金的种类很多，它们在汽车、船舶、飞机等制造业上以及在日常生活里的用途很广。

镁的主要用途也是制造各种轻合金。例如，镁跟铝、铜、 锡、锰、钛、铍等金属元素形成的多种合金（约含镁 80\%），密度只有 \qty{1.8}{g/cm^3} 左右，但硬度和强度都较大。
因此，镁合金也是汽车、飞机制造业的重要材料。

\subsection{化学性质}
镁和铝的原子容易失去最外层电子成为阳离子。
\begin{gather*} 
  \ce{ Mg - 2e -> Mg^2+} \\ 
  \ce{ Al - 3e -> Al^3+}
\end{gather*}

镁和铝都是比较活泼的金属，都是强还原剂，它们能跟非金属、酸等物质起反应。
铝还可以跟强碱溶液起反应。

\subsubsection{跟非金属的反应}
在常温下，镁和铝在空气里都能跟氧气起反应，生成一层致密而坚固的氧化物薄膜，从而使金属失去光泽。
但是，这层氧化物薄膜能够阻止金属的继续氧化，所以，镁和铝都有抗腐蚀的性能。

镁在空气里点燃可以燃烧，放出大量的热并发出耀眼的白光。
我们利用镁的这种性质来制造照明弹等。
\[ \ce{ 2Mg +O2 \xlongequal{\text{点燃}} 2MgO + \text{热量}} \]

铝跟氧气起反应没有镁那样容易，可是把铝粉或铝箔放在氧气里加热，铝也能燃烧，放出大量的热，同时发出耀眼的白光。
但是，在空气里，只是在高温下，铝才能发生这样剧烈的变化。
\[ \ce{ 4Al + 3O2 \xlongequal{\quad} 2Al2O3 + \text{热量} }\]

镁、铝除能跟氧气起反应外，还能跟其他非金属如硫、卤素等起反应。

\subsubsection{跟酸的反应}
镁和铝都能跟稀盐酸或稀硫酸起反应生成氢气。
在常温下，铝在浓硝酸里表面被钝化，生成坚固的氧化膜，可以阻止反应的继续进行。
因此，可以用铝制的容器来装运浓硝酸。

\subsubsection{跟碱的反应}
镁不能跟碱起反应，但铝能跟强碱溶液起反应，生成氢气和偏铝酸盐。
\[ \ce{2Al + 2NaOH +2H2O \xlongequal{\quad} 2NaAlO2 + 3H2 ^}\]

\subsubsection{跟某些氧化物的反应}
镁不仅能跟空气中的氧气起反应，而且能够跟二氧化碳起反应，夺取其中的氧，析出游离态的碳。

\begin{Experiment}*[righthand ratio=0.7]
如\cref{fig:6-3} 所示，点燃镁条，放入盛有二氧化碳的集气瓶里（瓶里装少量细砂），镁条剧烈燃烧，生成氧化镁，析出的碳附着在集气瓶内壁上。
\tcblower 
\begin{figurehere}
  \includegraphics{6-3.pdf}
  \caption{镁条在二氧化碳里燃烧}\label{fig:6-3}
\end{figurehere}
\end{Experiment}
反应的化学方程式如下：
\[ \ce{ 2Mg +CO2 \xlongequal{\text{燃烧}} 2MgO +C }\]

铝在一定条件下能够跟氧化铁发生氧化-还原反应。
% \begin{figure}
%   \begin{minipage}[b]{0.48\linewidth}\centering
%     \caption{镁条在二氧化碳里燃烧}\label{fig:6-3}
%   \end{minipage}
%   % \begin{minipage}[b]{0.48\linewidth}\centering
%   %   \includegraphics{6-4.pdf}
%   %   \caption{铝热反应实验装置}\label{fig:6-4}
%   % \end{minipage}
% \end{figure}

\begin{Experiment}*[righthand ratio=0.4]
用两张圆形滤纸分别折叠成漏斗状，套在一起使四周都叠上四层。
把内层的纸漏斗取出，在底部剪一小孔，用水润湿，再跟另一纸漏斗套在一起，架在铁圈上（如\cref{fig:6-4}），下面放置盛砂的蒸发皿。
把 \qty{5}{g} 干燥的氧化铁粉末和 \qty{2}{g} 铝粉混和均匀，放在纸漏斗里，上面加少量氯酸钾并在混和物中间插一根镁条，用小木条点燃镁条，观察发生的现象。
\tcblower 
\begin{figurehere}
  \includegraphics{6-4.pdf}
  \caption{铝热反应实验装置}\label{fig:6-4}
\end{figurehere}
\end{Experiment}

从实验我们看到，镁条剧烈燃烧，放出一定的热量，使氧化铁粉末和铝粉在较高的温度下发生了剧烈的反应，放出大量的热，发出耀眼的光芒。
我们还可以看到纸漏斗下部被烧穿，有熔融物落入砂中。
冷却后除去外层熔渣，观察，可以发现落下的是铁珠。
这个反应叫铝热反应。
在反应中温度可达 \qty{2000}{\celsius} 以上，生成了氧化铝和液态的铁。
\[ \ce{ 2Al + Fe2O3 \xlongequal{\quad} 2Fe + Al2O3 + \text{热量} } \]

通常把铝粉和氧化铁的混和物叫铝热剂。

\begin{Reading}[补充阅读]{}
铝热反应的原理可以应用在生产上，例如焊接钢轨。
往铝热剂里加入一定量的铁合金和铁钉屑，把这种混和物放到特制的坩埚里，点燃后,立即发生剧烈的复杂的化学反应\footnote{除发生铝热反应外，还有炼钢和造渣等反应。}，产生了高温的钢水和熔渣。
随即把高温钢水浇入扣在两根钢轨缝上的砂型中，这样可以把两根钢轨焊接起来（\cref{fig:6-5}）。
\begin{figurehere}
  \includegraphics{6-5.pdf}
  \caption{用铝热剂焊接钢轨}\label{fig:6-5}
\end{figurehere}
这种焊接不用电源，而且焊接的速度快，设备简易，适合在野外作业。
用这种焊接方法可以把比较短的钢轨（\qtyrange{12.5}{25}{m}）焊接成较长的钢轨（约 \qty{1000}{m}），这在建设铁道的无缝线路上有着重要的意义。
此外，在工业上，也用铝热剂的焊接法焊接大截面的钢材部件。
\end{Reading}

不仅用铝粉和氧化铁可以做铝热剂，用某些金属氧化物（如 \ce{V2O5}、\ce{Cr2O3}、\ce{MnO2} 等）代替氧化铁也可以做铝热剂。
当铝粉跟某些金属氧化物反应时，产生足够的热量，使被还原的金属在较高的温度下，呈熔融状态并跟形成的熔渣分离开来，从而获得较纯的金属。
在工业上常用这种方法冶炼难熔的金属，如钒、铬、锰等。

\begin{Practice}[习题]
\begin{question}
  \item 镁和铝都是较活泼的金属，为什么它们却不容易被腐蚀？
  \item 下列说法是否正确？为什么？
  \begin{tasks}[label=(\arabic*)]
    \task 因为二氧化碳是一种很好的灭火剂，所以在镁粉燃烧时，可以用二氧化碳来扑灭。
    \task 在化学反应里，镁原子能失去 2 个电子，而铝原子可以失去 3 个电子，所以铝的活泼性比镁的强。
  \end{tasks}
  \item 分别写出铝跟氢氧化钠、稀硫酸起反应的化学方程式，然后再把这两个化学方程式改写成离子方程式。
  \item 铝的下列用途主要是由它的哪些性质决定的？
  \begin{tasks}(3)
    \task 家用铝锅，
    \task 浓硝酸容器，
    \task 导线，
    \task 包装铝箔，
    \task 焊接钢轨。
  \end{tasks}
  \item 试比较第三周期元素的物理性质和化学性质，怎样利用这些性质说明元素周期表从左到右性质递变的规律？
  \item 写出钠、镁、铝分别跟足量的稀盐酸起反应的离子方程式，并计算：
  \begin{tasks}[label=(\arabic*)]
    \task 当金属各为 \qty{1}{g} 时，哪个反应放出的氢气多？
    \task 当金属各为 \qty{0.1}{mol} 时，哪个反应放出的氢气多？
  \end{tasks}
  \item 配平下列化学方程式，并指出哪些物质是氧化剂，哪些物质是还原剂？并计算生成 \qty{1}{mol} 各种金属，分别需要纯铝多少克？
  \begin{tasks}(2)
    \task $\ce{ V2O5  + Al }$；
    \task $\ce{ WO3   + Al }$；
    \task $\ce{ Cr2O3 + Al }$；
    \task $\ce{ Co3O4 + Al }$。
  \end{tasks}
\end{question}
\end{Practice}


\section{镁和铝的重要化合物\texorpdfstring{\quad}{ }铝的冶炼}
镁和铝在自然界的分布都很广。
铝是地壳里存在最多的金属元素，约占地壳总质量的 7.7\%。
镁的储量也很丰富，约占地壳总质量的 2.0\%。
由于它们都是较活泼的金属，因此都以化合态存在于自然界里。

含镁的主要矿物有光卤石、白云石、菱镁矿、硫酸镁、滑石和石棉等。
此外，海水中也含有大量的镁盐。

含铝的主要矿物有长石、云母、铝土矿和明矾石等。

下面简单介绍几种镁和铝的重要化合物。
\subsection{镁的重要化合物}
\subsubsection{氧化镁（\ce{MgO}）}
镁在空气里燃烧可以生成氧化镁。
工业上通常由煅烧菱镁矿（主要成分是 \ce{MgCO3}）来制取氧化镁：
\[ \ce{ MgCO3 \xlongequal{\text{煅烧}} MgO + CO2 ^ }\]

氧化镁是很轻的白色粉末,它的熔点高达 \qty{2800}{\celsius}，是优良的耐火材料，常用来制造耐火砖、耐火管和坩埚等。

氧化镁是一种碱性氧化物，可以跟水缓慢地反应生成氢氧化镁：
\[ \ce{ MgO + H2O \xlongequal{\quad} Mg(OH)2}\]

\subsubsection{氯化镁（\ce{MgCl2}）}
氯化镁是一种无色、苦味、易溶的晶体，极易吸收空气中的水分而潮解，粗制食盐在潮湿的空气中容易吸湿受潮，就是由于里面含有少量氯化镁杂质的缘故。

电解熔融的氯化镁，可以制得金属镁。
因此，氯化镁是制取镁的重要原料。

光卤石（\ce{KCl.MgCl2.6H2O}）中含有氯化镁，把光卤石溶解，提取出氯化钾以后，把溶液浓缩，就得到六水合氯化镁（\ce{MgCl2.6H2O}）晶体。

除了陆地上有大量含镁的矿物以外，海洋中也储存着大量的镁盐，海水中镁离子的平均含量为 0.13\%，仅次于氯和钠，居第三位。
由于海水总量大，因此海水是镁的主要来源之一。

下面简单介绍怎样从海水中提取镁。

这类工厂以大量海水为原料，都建在海边。
一般先把海滩上的贝壳煅烧成石灰，将石灰制成石灰乳。
把海水引入水渠，然后把石灰乳加到水中，由于氢氧化镁的溶解度很小，就生成氢氧化镁沉淀，经沉降、洗涤、过滤就可以得到氢氧化镁。 
这时，如果把氢氧化镁煅烧，可制得氧化镁。
如果在氢氧化镁中加入盐酸，则可制得氯化镁，经结晶、过滤、干燥，然后在电解槽中电解，可以得到金属镁和氯气，其中的氯气可制成盐酸循环使用。

除了上面介绍的方法以外，还可以从制取食盐剩下的苦卤中提取氯化镁。
苦卤里含有较多的氯化镁和氯化钠，其次是硫酸镁，氯化钾的含量则较少，此外，还含有少量溴化镁。
如果向苦卤里加入石灰，可生成氢氧化镁沉淀，再用上面介绍的方法制取氯化镁。
此外，还可以利用各种盐在水里溶解的性质的不同，通过改变溶液条件，使各种盐分别结晶出来。

\subsection{铝的重要化合物}
\subsubsection{氧化铝（\ce{Al2O3}）}
自然界存在的铝的矿物主要有铝土矿。
铝土矿又称矾土，它是由一水合氧化铝（\ce{Al2O3.H20}）、三水合氧化铝（\ce{Al2O3.3H20}）以及少量氧化铁和石英等杂质组成的。
铝土矿可用来提取纯氧化铝。
氧化铝是一种白色难熔的物质，是冶炼铝的原料。
氧化铝也是一种比较好的耐火材料，它可以用来制造耐火坩埚、耐火管和耐高温的实验仪器等。

\begin{Reading}[补充阅读]{}
刚玉是天然产的无色氧化铝晶体。
刚玉的硬度很大，仅次于金刚石。
因此，它常被用来制成砂轮、研磨纸或研磨石等，用以加工光学仪器和某些金属制品。
通常所说的蓝宝石和红宝石是混有少量不同氧化物杂质的刚玉。
它们可以用做精密仪器和手表的轴承。
\end{Reading}

氧化铝是典型的两性氧化物，新制备的氧化铝既能跟酸起反应生成铝盐，又能跟碱起反应，生成偏铝酸盐。
\begin{gather*}
  \ce{  Al2O3 + 6H+ \xlongequal{\quad} 2Al^3+ + 3H2O}\\
  \ce{  Al2O3 + 2OH- \xlongequal{\quad} 2AlO2- + H2O}
\end{gather*}

氧化铝不溶于水，因此，不能用它来制备氢氧化铝。

\subsubsection{氢氧化铝（\ce{Al(OH)3}）}
氢氧化铝是不溶于水的白色胶状物质。
氢氧化铝能凝聚水中悬浮物，又有吸附色素的性能。
在实验室里我们可以用铝盐溶液跟氨水的反应来制备氢氧化铝。

\begin{Experiment}
  把 \qty{3}{mol} \qty{0.5}{M} 硫酸铝溶液注入试管里，向硫酸铝溶液里滴加氨水，生成蓬松的白色胶状氢氧化铝沉淀，继续滴加氨水，直到不再产生沉淀为止。
过滤，用蒸馏水冲洗沉淀，得到较纯净的氢氧化铝（大部分留下做下面的实验）。
取少量氢氧化铝沉淀放在蒸发皿里，加热。
观察氢氧化铝的分解。
\end{Experiment}
上述实验的反应可表示如下：
\begin{gather*}
  \ce{ Al^3+ +3NH3.H2O \xlongequal{\quad} Al(OH)3 v + 3NH4+ }\\
  \ce{ 2Al(OH)3 \xlongequal{\triangle} Al2O3 +3H2O }
\end{gather*}

\begin{Experiment}
把上面实验制得的氢氧化铝沉淀分装在两个试管里，往一个试管里滴加 \qty{2}{M} 盐酸；往另一个试管滴加 \qty{2}{M} 氢氧化钠溶液。观察两个试管里发生的现象。
\end{Experiment}

通过上述实验我们看到，氢氧化铝在酸或强碱的溶液里都能溶解。
可见它既能跟酸起反应，又能跟强碱起反应。这两个反应可以分别表示如下：
\begin{gather*}
  \ce{ Al(OH)3 + 3H+ \xlongequal{\quad} Al^3+ + 3H2O }\\
  \ce{ Al(OH)3 + OH- \xlongequal{\quad} AlO2- + 2H2O }
\end{gather*}

实验证明，氢氧化铝是典型的两性氢氧化物。

为什么氢氧化铝会具有两性呢？
我们可以用酸、碱的电离以及平衡移动原理来简单地加以解释。

氢氧化铝的电离方程式可表示如下：
\[ \ce{ H2O + AlO2- + H+ <=> Al(OH)3 <=> Al^3+ + 3OH- } \]

氢氧化铝是一种弱电解质，它的碱性和酸性都很弱。
也就是说，它电离时生成的 \ce{H+} 和 \ce{OH-} 都很少。
那么，它为什么既能跟酸起反应，又能跟碱起反应呢？ 
我们可以用平衡移动原理加以说明。

当往 \ce{Al(OH)3} 里加酸时，\ce{H+} 立即跟溶液里少量的 \ce{OH-} 起反应而生成水，这就会使氢氧化铝按碱式电离，使平衡向右移动，从而使 \ce{Al(OH)3} 不断地溶解。
相反地,当往 \ce{Al(OH)3} 里加碱时，\ce{OH-} 立即跟溶液里少量的 \ce{H+} 起反应而生成水，这样就会使氢氧化铝按酸式电离，使平衡向左移动，同样 \ce{Al(OH)3} 也就不断地溶解了。

\subsubsection{硫酸铝钾（\ce{KAl(SO4)2}）}
硫酸铝钾是由两种不同的金属离子和一种酸根离子组成的盐，象这样的盐，叫做\Concept{复盐}。它电离时可解离出两种金属的阳离子。
\[ \ce{ KAl(SO4)2 \xlongequal{\quad} K+ +Al^3+ + 2SO4^2- } \]

十二水合硫酸铝钾（\ce{KAl(SO4)2.12H2O}）的俗名是明矾，明矾通常也使用 \ce{K2SO4.Al2(SO4)3.24H2O} 表示。明矾是无色晶体，易溶于水，跟水发生水解反应，它的水溶液呈酸性。
\[ \ce{ Al^3+ +3H2O <=> Al(OH)3\text{ （胶体）} + 3H+ } \]

明矾水解所产生的 \ce{Al(OH)3} 胶体吸附能力很强，可以吸附水里悬浮的杂质，并形成沉淀，使水澄清。
所以，明矾是一种较好的净水剂。

\subsection{铝的冶炼}
我们已经知道，铝是地壳里含量最多的金属，它有许多优良的性质，被普遍应用在各个领域里。
铝虽然储量多，用途广，但人们发现铝的历史并不长，这是因为铝的化学性质活泼，不易还原，因此从矿石中炼铝就比较困难。
一百多年以前，人们曾用金属钠还原氧化铝制得金属铝，由于当时钠很贵，因而铝比钠更贵，价格跟黄金不相上下，因此当时被认为是一种稀罕的贵金属，并且被制成最珍贵的奖品、礼品。
直到十九世纪末，人们发明了生产铝的新方法——电解冰晶石\footnote{冰晶石又叫六氟合铝酸钠，它以晶体的光泽近于冰而得名，通常是一种无色或呈灰色的物质。}和氧化铝的熔融混和物，才能够大规模冶炼金属铝。
自此以后，铝的价格大大下降，铝才真正获得了广泛的应用。
现在，工业上就是采用这个方法制铝的。

纯净氧化铝的熔点很高（约为 \qty{2045}{\celsius}）很难熔化。
用熔化的冰晶石（\ce{Na3AlF6}）作熔剂，可使氧化铝在 \qty{1000}{\celsius} 左右溶解在液态的冰品石里，成为冰晶石和氧化铝的熔融体，然后进行电解。

冶炼铝的主要设备是电解槽，如\cref{fig:6-6} 所示。
电解槽呈长方形，外面是钢壳，内衬耐火砖作为绝缘和保温层。
用碳块作槽池，底部碳块跟它下部的钢导电棒相连结作阴极。
用垂直于电解槽带导电杆的两排碳块作阳极。
中间是冰晶石—氧化铝熔融体。
\begin{figure}
  \includegraphics{6-6.pdf}
  \caption{冶炼铝的电解槽}\label{fig:6-6}
\end{figure}

在直流电通过冰晶石—氧化铝的熔融体时，发生了很复杂的反应。可简略地用下式表示两极的反应：

在阴极
\[ \ce{ 4Al^3+ +12\mathrm{e} \xlongequal{\quad} 4Al } \]

在阳极
\[ \ce{6O^2- -12\mathrm{e} \xlongequal{\quad} 3O2 } \]

总的反应
\[ \ce{2Al2O3 \xlongequal{\text{电解}} 4Al + 3O2 ^ }\]

在阳极生成的氧气立即跟碳极起反应，生成二氧化碳。
\[ \ce{C + O2 \xlongequal{\quad} CO2 ^} \]
因此，在电解过程里，氧化铝和碳块需要定期补充。

由于温度在 \qty{1000}{\celsius} 左右，电解生成的铝在电解槽里呈液态铝析出。
冰晶石—氧化铝熔融体的密度小于液态铝的密度，因而液态铝就积存在槽底，可以定期汲出。

\begin{Practice}[习题]
\begin{question}
  \item 试用化学方程式表示下列各步反应：
  \begin{tasks}
    \task \ce{Mg}$\to$\ce{MgO}$\to$\ce{MgCl2}$\to$\ce{MgSO4}$\to$\ce{Mg(OH)2}$\to$\ce{MgCl2}$\to$\ce{Mg}
    \task \ce{Al}$\to$\ce{Al2O3}$\to$\ce{AlCl3}$\to$\ce{Al2(SO4)3}$\to$\ce{Al(OH)3}$\to$\ce{NaAlO2}
  \end{tasks}
  \item 铝锅表面既然蒙有一层有保护作用的氧化膜，为什么还不宜用碱永洗涤？为什么不宜用来蒸煮酸的食物？
  \item 有一块镁铝合金，把它溶于盐酸里，再加过量的氢氧化钠溶液，问镁、铝各以什么形式存在？写出反应的化学方程式。
  \item 有两位同学分别设计了制取 \ce{Al2S3} 的方案，请你判断哪位同学的方案正确，并说明理由。
  \begin{tasks}
    \task 用铝粉和硫粉在高温下反应制取。
    \task 用硫化钠溶液和氯化铝溶液反应制取。
  \end{tasks}
  \item 有一包白色粉末，已知它是氯化镁、氢氧化镁或碳酸镁中的一种，试设计一个实验，判断它是哪一种物质。
  \item 某铝土矿里含有氧化铝（\ce{Al2O3}）88\%，在理论上，生产 \qty{10}{t} 金属铝，需要这种矿石多少吨？
\end{question}
\end{Practice}

\section{硬水及其软化}
\subsection{水的硬度}
水是日常生活和王农业生产不可缺少的物质。
水质的好坏对生产和生活影响很大。
天然水跟空气、岩石和土壤等长期接触，溶解了许多杂质: 如无机盐类、某些可溶性有机物以及气体等。
水里溶解的无机盐有钙和镁的酸式碳酸盐、碳酸盐、氯化物、硫酸盐、硝酸盐，等等。
也就是说，天然水里一般地含有 \ce{Ca^2+}、\ce{Mg^2+} 等阳离子和 \ce{HCO3-}、\ce{CO3^2-}、\ce{Cl-}、\ce{SO4^2+} 和 \ce{NO3-} 等阴离子。
各地的天然水含有这些离子的种类和数量有所不同。

有的天然水里含 \ce{Ca^2+}、\ce{Mg^2+} 比较多，在这种水里加入肥皂水，就会发生沉淀。

\begin{Experiment}\label{expr:6-5}
  用两个试管分别取蒸馏水、天然水各 \qtyrange{5}{6}{ml}，各注入少量肥皂水，振荡。观察发生的现象。
\end{Experiment}

我们可以看到盛有蒸馏水的试管里泡沫很多，没有沉淀产生。
在盛有天然水的试管里泡沫较少，并出现絮状沉淀。
这是因为天然水里含有 \ce{Ca^2+} 或 \ce{Mg^2+}，肥皂跟 \ce{Ca^2+} 或 \ce{Mg^2+} 起反应以后，生成了不溶于水的物质 \footnote{肥皂的主要成分是硬脂酸钠（\ce{Na(C12H35O2)}），肥皂跟 \ce{Ca^2+}（或 \ce{Mg^2+}）起反应生成的 \ce{Ca(C12H35O2)2}（或 \ce{Mg(C12H35O2)2}）是不溶于水的物质。}。各地天然水所含有的 各种离子的量不同,一般地说,地下水、泉水中含 \ce{Ca^2+}、\ce{Mg^2+} 较多，而河水、湖水、雨水中含得少一些。

通常把溶有较多量的 \ce{Ca^2+} 和 \ce{Mg^2+} 的水，叫做\Concept{硬水}；只溶有少量或不含 \ce{Ca^2+} 或 \ce{Mg^2+} 的水，叫做\Concept{软水}。

如果水的硬度\footnote{水的硬度常用一种规定的标准来衡量。通常把一升水里含有 \qty{10}{mg} \ce{CaO}（或相当于 \qty{10}{mg} \ce{CaO} 称为 1 度（以 \ang{1} 表示）。水的硬度在 \ang{8} 以下的称为软水; 在 \ang{8} 以上的称为硬水。硬度大于 \ang{30} 的是最硬水。}是由碳酸氢钙或碳酸氢镁所引起的，这种硬度叫做\Concept{暂时硬度}。
具有暂时硬度的水经过煮沸以后，水里所含的碳酸氢钙就分解而生成不溶性的碳酸钙:
\[ \ce{Ca(HCO3)2 \xlongequal{\triangle} CaCO3 v + CO2 ^ + H2O} \]
水里所含的碳酸氢镁先形成难溶于水的碳酸镁沉淀。
\[ \ce{Mg(HCO3)2 \xlongequal{\triangle} MgCO3 v + CO2 ^ + H2O} \]
碳酸镁虽然难溶于水，但它仍有少量能溶解在水里（在 \qty{18}{\celsius} 时，它的溶解度为 \qty{0.011}{g}，相当于 \qty{110}{mg/l}）。
因此，当继续加热煮沸时，碳酸镁就发生水解反应，生成了更难溶的氢氧化镁（在 \qty{18}{\celsius} 时，它的溶解度为 \qty{0.00084}{g}，相当于 \qty{8.4}{mg/l}）。
这样,水里溶解的 \ce{Ca^2+} 和 \ce{Mg^2+} 就成为碳酸钙和氢氧化镁的沉淀，从水里析出，水的硬度就可以降低，从而使硬度高的水得到软化。

如果水的硬度是由钙和镁的硫酸盐或氯化物等所引起的，这种硬度叫做\Concept{永久硬度}。
永久硬度不能用加热的方法软化。
天然水大多同时具有暂时硬度和永久硬度，因此，一般所说水的硬度是泛指上述两种硬度的总和。


\subsection{硬水的软化}
水的硬度高对生活和生产都有危害。
洗涤用水如果硬度太高，不仅浪费肥皂，而且衣物也不易洗干净。
长期饮用硬度很高或硬度过低的水，都不利于人体的健康。
锅炉用水硬度太高，特别是暂时硬度太高，十分有害。
因为经过长期烧煮后，水里的钙盐和镁盐会在锅炉内结成锅垢，使锅炉内的金属管道的导热能力大大降低，这不但浪费燃料，而且会使管道局部过热，当超过金属允许的温度时，锅炉管道将变形或损毁，严重时会引起锅炉爆炸事故。
很多工业部门，如纺织、印染、 造纸、化工等，都要求用软水。
因此，对天然水进行处理，以降低或消除它的硬度是很重要的。

软化硬水的方法通常有药剂软化法和离子交换法等。

\subsubsection{药剂软化法}
药剂软化法是在水里加入适当的药剂。
使溶解在水里的钙、镁盐转化为几乎不溶于水的物质，生成沉淀从水里析出，以达到软化水的目的。
常用的药剂有石灰、纯碱或磷酸三钠等。
根据对水质的要求，可以用一种药剂，也可以同时用几种药剂。
例如，石灰纯碱软化水的方法，就是用石灰和纯碱（\ce{Na2CO3}）跟溶解在水里的 \ce{Ca^2+}、\ce{Mg^2+} 生成沉淀，从而使硬水得到软化。

石灰跟水里的碳酸氢钙和碳酸氢镁起反应。
\begin{gather*} 
  \ce{Ca(HCO3)2 + Ca(OH)2 \xlongequal{\quad} 2CaCO3 v + 2H2O }\\
  \ce{Mg(HCO3)2 + 2Ca(OH)2 \xlongequal{\quad} 2CaCO3 v + Mg(OH)2 v + 2H2O }
\end{gather*}
石灰也可以跟水里的硫酸镁或氯化镁起反应，生成氢氧化镁沉淀和硫酸钙或氯化钙。
这样把由 \ce{Mg^2+} 所形成的水的硬度转化为由 \ce{Ca^2+} 所形成的水的硬度。
\[ \ce{ MgSO4 + Ca(OH)2 \xlongequal{\quad} Mg(OH)2 v + CaSO4 } \]
加入纯碱，可以除去由 \ce{Ca^2+} 所形成的水的硬度。
\[ \ce{ CaSO4 + Na2CO3 \xlongequal{\quad} CaCO3 v + Na2SO4 }\]

\subsubsection{离子交换法}
离子交换法是用离子交换剂\footnote{离子交换剂包括天然或人造沸石、磺化煤和离子交换树脂等物质。}软化水的一种现代的方法。

在工业上常用磺化煤\footnote{磺化煤是烟煤、褐煤用发烟硫酸或浓硫酸处理后的产物。为了简便起见，可以用 \ce{NaR} 表示。}（\ce{NaR}）做离子交换剂。
磺化煤是黑色颗粒状物质，它不溶于酸和碱。
这种物质的阳离子会跟溶液里其他物质的阳离子发生离子交换作用。
\par\medskip\noindent
\begin{minipage}{0.6\linewidth}\parindent2em
如\cref{fig:6-7} 所示，在离子交换柱里装有磺化煤。
把硬水从离子交换柱的上口注入，使硬水慢慢地流经磺化煤。
当硬水通过磺化煤时，硬水里的 \ce{Ca^2+} 和 \ce{Mg^2+} 跟磺化煤的 \ce{Na+} 起离子交换作用，因而使硬水得到软化，反应可以表示如下：
\begin{gather*}
  \ce{2NaR + Ca^2+ \xlongequal{\quad} CaR2 +2Na+}\\
  \ce{2NaR + Mg^2+ \xlongequal{\quad} MgR2 +2Na+}
\end{gather*}
\end{minipage}\hfill
\begin{minipage}{0.35\linewidth}
\begin{figurehere}
  \includegraphics{6-7.pdf}
  \caption{用离子交换剂软化硬水}\label{fig:6-7}
\end{figurehere}
\end{minipage}\par\medskip


\begin{Experiment}
  把\cref{expr:6-5} 所用的天然水经离子交换剂软化以后，用试管取 \qtyrange{5}{6}{ml}，注入少量肥皂水，振荡。
\end{Experiment}
通过实验，我们可以看到，硬水软化以后，加入肥皂水振荡时，泡沫很多，没有沉淀产生。

磺化煤的 \ce{Na+} 全部为 \ce{Ca^2+} 和 \ce{Mg^2+} 所代替后，磺化煤就失去了软化硬水的能力。
但是用 8\%～10\% 食盐溶液浸泡，\ce{CaR2} 和 \ce{MgR2} 跟 \ce{Na+} 就又会发生交换作用，重新生成 \ce{NaR}，从而又恢复了磺化煤软化硬水的能力，这个过程叫做再生。可以表示如下：
\[ \ce{ CaR2 +2Na+ \xlongequal{\quad} 2NaR +Ca^2+ }\]

这种软化硬水的方法比药剂软化法具有质量高，设备简单，占用面积小，操作方便等优点，因此，目前使用比较普遍。

\begin{Practice}[习题]
\begin{question}
  \item 某井水里既含有硫酸钙，又含有碳酸氢钙，怎样把它软化？ 写出有关反应的化学方程式。
  \item 以石灰纯碱软化水的方法为例，说明药剂软化水的化学原理是什么？写出有关的化学方程式。用这种方法软化的水，为什么仍有一定的硬度？
  \item 试回答下列问题：
  \begin{tasks}[label=(\arabic*)]
    \task 明矾为什么有净水作用？用明矾处理硬水以后，硬水是否变为软水？
    \task 用硬水洗衣服会增加肥皂的消耗量，为什么？
    \task 除去了天然水中的 \ce{HCO3-}、\ce{SO4^2-}、\ce{Cl-}，是否就把硬水软化了？
  \end{tasks}
\end{question}
\end{Practice}

\section*{内容提要}
\addcontentsline{toc}{section}{内容提要}\setcounter{subsection}{0}
\subsection*{一、金属键}
金属晶体里的金属离子跟自由电子之间存在着较强的作用，这种作用叫做金属键。
金属键没有方向性。
金属键也是化学键的一种。

金属的一些共同的性质，如金属光泽、易导电、易导热、有延展性等等，都可以由金属键得到一定的解释。

\subsection*{二、镁和铝的性质}\setcounter{subsubsection}{0}
\subsubsection{物理性质}
镁和铝的密度小，有延展性，铝易导电。
镁和铝能跟许多金属或某些非金属形成合金。

合金就是两种或两种以上的金属（或金属跟非金属）熔合而成的具有金属特性的物质。
\subsubsection{化学性质}
\begin{enumerate}
  \item 镁和铝能跟许多非金属、稀盐酸（或稀硫酸）起反应。 铝在浓硝酸里被钝化。
  \item 镁不能跟碱起反应；铝能跟强碱溶液起反应生成氢气。
  \item 镁和铝能跟某些氧化物起反应：
  
  镁能跟二氧化碳起反应。

  铝能跟氧化铁起反应，生成氧化铝和液态铁，并产生大量的热。
  铝粉跟氧化铁的混和物叫做铝热剂。

  用某些金属氧化物（如 \ce{V2O5}、\ce{Cr2O3}等）代替氧化铁跟铝起反应，能制取某些金属（如钒、铬等）。
\end{enumerate}

\subsection*{三、镁和铝的重要的化合物}
镁的重要化合物有氧化镁和氯化镁等。
氯化镁是电解法制镁的原料。

铝的重要化合物有氧化铝、氢氧化铝和明矾等。
氧化铝和氢氧化铝都是两性化合物。
氢氧化铝跟酸或跟强碱起反应，可用离子方程式表示如下：
\begin{align*}
  \ce{ Al(OH)3 + 3H+ \xlongequal{\quad} Al3+ +3H2O }\\
  \ce{ Al(OH)3 + OH- \xlongequal{\quad} AlO2- +2H2O }
\end{align*}

氢氧化铝的电离方程式可表示如下：
\[ \ce{H+ +AlO2- +H2O <=> Al(OH)3 <=> Al3+ +3OH-} \]

明矾（\ce{KAl(SO4)2.12H2O}）是一种复盐。明矾容易起水解反应生成 \ce{Al(OH)3} 胶体，它常用做净水剂。

\subsection*{四、铝的冶炼}
铝的冶炼是把纯氧化铝溶解于熔融的冰晶石（\ce{Na3AlF6}）里作电解质，在电解槽里进行电解。
电解时的化学反应很复杂，可简略地用化学方程式表示如下：
\[ \ce{2Al2O3 \xlongequal{\text{电解}} 4Al + 3O2 ^}\]

\subsection*{五、硬水及其软化}
\begin{enumerate}[1.]
  \item 含有较多 \ce{Ca^2+} 和 \ce{Mg^2+} 的水叫做硬水。含有较少（或不含有）\ce{Ca^2+} 和 \ce{Mg^2+} 的水叫做软水。
  
  由钙和镁的酸式碳酸盐所引起的水的硬度叫暂时硬度。

  由钙和镁的硫酸盐或氯化物所引起的水的硬度叫永久硬度。

  水的硬度一般地是暂时硬度和永久硬度的总和。
  \item 硬水的软化
  
  硬水对生活和生产都有危害，一般常用的软化硬水的方法有药剂软化法和离子交换法等。
\end{enumerate}

\begin{Review}
\begin{question}
  \item 镁在氮气里燃烧生成氮化镁（\ce{Mg3N2}），氮化镁跟水起反应生成氢氧化镁和氮气。写出这两个反应的化学方程式并计算：
  \begin{tasks}[label=(\arabic*)]
    \task \qty{1.2}{g} 镁完全跟氮气起反应，问需要氮气多少毫升？（在 \qty{17}{\celsius} 和压强为 \qty{750}{mmHg}）
    \task 在相同温度和压强下，可得到氮气多少毫升？
  \end{tasks}
  \item 在有固体 \ce{Mg(OH)2} 存在的饱和水溶液中，达到如下平衡：
  \[ \ce{ Mg(OH)2 \text{ （固）} <=> Mg^2+ + 2OH-}\]
  向其中分别加入下列物质后，\ce{Mg(OH)2} 固体的量发生什么变化？并简单加以说明。
  \begin{tasks}(6)
    \task 水，
    \task*(2) \ce{Mg(OH)2}（固），
    \task*(2) \ce{NaOH}（固），
    \task 盐酸。
  \end{tasks}
  \item 氢氧化铝的溶解度很小，所以由它电离而生成的离子当然也很少。为什么不管加入酸或强碱的溶液都能使氢氧化铝全部溶解而得到完全透明的溶液？试根据平衡移动的原理加以解释。
  \item 用含有 \ce{Ag+}、\ce{Al^3+}、\ce{Ca^2+}、\ce{Mg^2+} 等四种阳离子的溶液，做如下的实验，试填写下列的空格:
  \begin{tasks}[label=(\arabic*)]
    \task\label{tsk:experiment1} 往这种溶液里加入稀盐酸生成\CJKunderline[hidden]{\ce{AgCl}} 沉淀。
    \task\label{tsk:experiment2} 把实验 \ref{tsk:experiment1} 所得的液体过滤，往滤液加入氨水，使溶液呈碱性，又产生沉淀，沉淀是\CJKunderline[hidden]{\ce{AgCl}} 。
    \task\label{tsk:experiment3} 把实验 \ref{tsk:experiment2} 所得的液体过滤，往滤液里加入碳酸钠，又会产生 \CJKunderline[hidden]{\ce{AgCl}} 沉淀。
    \task\label{tsk:experiment4} 把实验 \ref{tsk:experiment2} 所得的沉淀中加入氢氧化钠溶液，直到过量，这时部分沉淀溶解生成 \CJKunderline[hidden]{\ce{AgCl}}，不溶解的是 \CJKunderline[hidden]{\ce{AgCl}}。
    \task 写出上述各化学反应的离子方程式。
  \end{tasks}
  \item 现有 $0.003M$ \ce{AgNO3} 溶液 \qty{600}{ml}，加入 $0.002M$ \ce{NaCl} 溶液 \qty{400}{ml}，起反应以后，在这种溶液里最多的是下列哪一种离子？
  \begin{tasks}(4)
    \task \ce{NO3-}，
    \task \ce{Ag+}，
    \task \ce{Cl-}，
    \task \ce{Na+}。
  \end{tasks}
  \item 电解食盐水时使用的食盐含有如下杂质：泥砂、\ce{CaCl2}、\ce{MgSO4}、\ce{CaSO4} 等，如何除去这些杂质？简述过程，写出有关反应的离子方程式。
\end{question}
\end{Review}

\begin{Exercise}*[总复习题]
\begin{question}
  \item 写出你所知道的，周期表前三周期的元素所形成的单质各属于哪种类型的晶体。
  \item 下列七种物质中，由非极性键形成的非极性分子是\CJKunderline*[hidden]{\hspace{1em}A、 G\hspace{1em}}，由极性键形成极性分子的是\CJKunderline[hidden]{\hspace{1em}B\hspace{1em}}，由极性键形成非极性分子的是\CJKunderline[hidden]{\hspace{1em}F\hspace{1em}}，只由离子键形成的离子晶体是\CJKunderline[hidden]{\hspace{1em}D\hspace{1em}}，既有离子键又有共价键的离子晶体是\CJKunderline*[hidden]{\hspace{1em}C、E\hspace{1em}}，既有离子键又有共价键、配位键的离子晶体是\CJKunderline[hidden]{\hspace{1em}C\hspace{1em}}。
  \begin{tasks}[label=\Alph*.](4)
    \task \ce{N2}
    \task \ce{NH3}
    \task \ce{NH4Cl}
    \task \ce{BaCl2}
    \task \ce{Ca(OH)2}
    \task \ce{CCl4}
    \task \ce{Cl2}
  \end{tasks}
  \item 根据下列物质的物理性质，推测它们在固态时形成哪类晶体。
  \begin{tasks}
    \task \ce{NaOH}\quad 熔点 \qty{318.4}{\celsius}，沸点 \qty{1390}{\celsius}，易溶于水，熔融时能导电。
    \task \ce{SO2}\quad 熔点 \qty{-72.7}{\celsius}，沸点 \qty{-10.08}{\celsius}，易溶于水。
    \task \ce{B}\quad 熔点 \qty{2300}{\celsius}，沸点 \qty{2550}{\celsius}，硬度大。
    \task \ce{CH4}\quad 熔点 \qty{-182.5}{\celsius}，沸点 \qty{-164}{\celsius}，易溶于有机溶剂。
  \end{tasks}
  \item 下列说法是不是正确，为什么？
  \begin{tasks}
    \task 红磷在空气中能自燃，又有毒性。
    \task 五氧化二磷易吸水，可以作为氨的干燥剂。
    \task 由于磷酸是一种二元酸，它能形成正盐和酸式盐。
    \task \ce{Ca(H2PO4)2} 比 \ce{CaHPO4} 和 \ce{Ca3(PO4)2} 易溶于水。
  \end{tasks}
  \item 怎样鉴别以下几种化肥: 硫酸铵、碳酸氢铵、氯化钾。
  \item 自然界和工业上有哪些氮的固定过程？为什么合成氨对工农业生产和国防十分重要？
  \item 氨的催化氧化反应 \ce{4NH3 + 5O2 ->[\ce{Pt}][\qty{800}{\celsius}] 4NO + 6H2O } 在一定条件下达到平衡状态。
  \begin{tasks}
    \task 这时，\CJKunderline[hidden]{\hspace*{1em}正反应\hspace*{1em}}和\CJKunderline[hidden]{\hspace*{1em}逆反应\hspace*{1em}}的反应速度相等，\CJKunderline[hidden]{\hspace*{1em}反应速度\hspace*{1em}}保持不变，而\CJKunderline*[hidden]{\hspace*{1em}正、逆反应\hspace*{1em}}都仍在继续进行。因此，化学平衡是一种\CJKunderline[hidden]{\hspace*{1em}动态\hspace*{1em}}平衡。
    \task \ce{Pt} 在这个反应中起\CJKunderline[hidden]{\hspace*{1em}催化剂\hspace*{1em}}作用，它能降低\CJKunderline[hidden]{\hspace*{1em}反应条件\hspace*{1em}}，加快\CJKunderline[hidden]{\hspace*{1em}化学反应\hspace*{1em}}的进行，但不能使\CJKunderline[hidden]{\hspace*{1em}反应平衡\hspace*{1em}}移动。
    \task 如果温度降低会使上述平衡向正反应方向移动，那么，正反应是\CJKunderline[hidden]{\hspace*{1em}放\hspace*{1em}}热反应。如果每摩尔氨氧化的反应热是 \qty{54.2}{kCal}，那么，这个反应的热化学方程式为\CJKunderline[hidden]{\hspace*{15em}}。
    \task 在上述平衡混和物中，如增大氧气的浓度，那么，上述平衡将向\CJKunderline[hidden]{\hspace*{1em}正\hspace*{1em}}方向移动。
  \end{tasks}
  \item 回答以下问题：
  \begin{tasks}
    \task 既然催化剂不能影响化学平衡的移动，为什么在化工生产上常常使用催化剂？
    \task 既然浓度的变化不能使平衡常数发生改变，为什么在反应中往往要增大反应物浓度？
    \task 既然合成氨是个放热反应，升高温度不利于反应向生成氨的方向进行，为什么合成氨厂常采用 \qty{500}{\celsius} 左右的反应温度？
  \end{tasks}
  \item 在某温度时，\ce{CO + H2O <=> CO2 + H2} 达到平衡后，$K=1$。\ce{CO} 的起始浓度为 \qty{0.01}{mol/L}，\ce{H2O}（气）的起始浓度为 \qty{0.03}{mol/L}，求达到平衡时，四种物质的浓度和 \ce{CO} 的转化率。
  \item 在某温度时，在容积为 \qty{4}{L} 的密闭器中装入 \qty{0.1}{mol} \ce{SO2} 和 \qty{0.05}{mol} 氧气，当达到平衡时，得到 \qty{0.06}{mol} \ce{SO3}，求该温度下反应的平衡常数和 \ce{SO2} 的转化率。
  \item 在某温度时，把 \qty{0.64}{mol} 的氢气和 \qty{0.64}{mol} 的碘蒸气通入 \qty{10}{L} 的容器里，当反应达到平衡时，生成 \qty{1}{mol} 碘化氢气体。求在这个温度下反应的平衡常数。
  \item 二氧化硅不能溶于水，怎样以二氧化硅为原料来制备硅酸？怎样以二氧化硅为原料来制取硅？
  \item 解释下列现象：
  \begin{tasks}
    \task 把鸡蛋白（液溶胶）倒入热锅里，立即凝结成固体。
    \task 在豆浆里加入蔗糖，除增加甜味外，并无什么现象发生，而加入盐卤（主要成分是 \ce{MgCl2}）或石膏（\ce{CaSO4.2H2O}），却发生凝聚作用，为什么？
    \task 用明矾（\ce{KAl(SO4)2.12H2O}）净水时，水中会出现 \ce{Al(OH)3} 胶体微粒。这种微粒是怎样形成的？它对清除水中的悬浮杂质起什么作用？
  \end{tasks}
  \item 电解精炼铜时，以粗铜片（内含 \ce{Ag}、\ce{Zn} 等杂质）作为阳极，以纯铜片作为阴极，以 \ce{CuSO4} 溶液为电解液，进行电解。回答下列问题:
  \begin{tasks}
    \task 写出阴极、阳极发生的化学反应。
    \task 为什么阴极不析出 \ce{Zn}？
    \task 阴极上析出的铜跟作为阳极的粗铜片在成分上有什么不同？
    \task 为什么在阳极附近沉积的阳极泥里可提取银？
  \end{tasks}
  \item 下列说法有没有错误，说明理由。
  \begin{tasks}
    \task 因为盐酸是强酸而醋酸是弱酸，如用相同当量浓度的盐酸和醋酸去中和同体积、同浓度的氢氧化钠溶液时，需用醋酸的量比盐酸多。
    \task 溶液的 pH 值减小一个单位，溶液的 \ce{[H+]} 增大 10 倍,溶液的 pH 值减小 2 个单位，溶液的 \ce{[H+]} 增大 20 倍。
    \task 因为强酸强碱盐不水解，它的水溶液呈中性，因此，凡是其水溶液呈中性的盐都是不水解的。
  \end{tasks}
  \item 从\cref{fig:5-2} 来看，表上所列 7 种弱酸中，哪种酸性相对较强，哪种酸性相对较弱？
  \item 已知常温时醋酸的电离常数为 \num{1.75e-5}，求 $1M$ 醋酸、$0.5M$ 醋酸、$0.1M$ 醋酸的电离度。
  \item 计算下列溶液的 pH 值
  \begin{tasks}(2)
    \task $0.1M$ \ce{HCl}
    \task $0.05M$ \ce{NaOH}
    \task $0.2M$ \ce{CH3COOH}
    \task $0.001M$ \ce{NH3.H2O}
  \end{tasks}
  \item 在 \qty{10}{ml} $0.5M$ \ce{HCl} 里加入 \qty{30}{ml} $0.1M$ \ce{NaOH}，求混和溶液的 pH 值。
  \item 等质量的镁和铝分别跟足量的稀 \ce{H2SO4} 完全反应，放出的氢气是不是一样多，为什么？等物质的量的镁和铝分别跟足量的稀 \ce{H2SO4} 完全反应，放出的氢气的量是不是一样多，为什么？
\end{question}
\end{Exercise}